%% ---------------------------------------------------------------------
%% 02 - ĐỐI TƯỢNG VÀ PHƯƠNG PHÁP NGHIÊN CỨU
%% Mô tả nền tảng thử nghiệm, dữ liệu, giả định và thuật toán.
%% Mục con dùng \subsection{...} (không đánh số, 10 pt sans bold).
%% ---------------------------------------------------------------------

\section{Đối tượng và phương pháp nghiên cứu}

\subsection{Nền tảng thử nghiệm và giả định}

Nền tảng thử nghiệm là robot vi sai \textit{Hermes} của nhóm tác giả, được
trang bị LiDAR 2D quét $270^\circ$ với tầm quan sát danh định $30$~m và một
máy tính nhúng chạy ROS~2. Môi trường thử nghiệm là hành lang văn phòng có
vật cản tĩnh và người đi lại, được ghi lại thành bản đồ lưới độ phân giải
$0{,}05$~m.

Các giả định: (i)~vị trí robot luôn xác định được nhờ SLAM; (ii)~vật cản
động chuyển động với vận tốc không vượt quá $1{,}5$~m/s; (iii)~bán kính an
toàn của robot là $0{,}35$~m.

\subsection{Biểu diễn hình học bằng Voronoi có trọng số}

Xét tập vật cản dạng đĩa $\{\mathbf{p}_i, r_i\}_{i=1}^{n}$ trong mặt phẳng.
Khoảng cách có trọng số từ điểm $\mathbf{q}$ tới vật cản thứ $i$ được định
nghĩa như trong phương trình~\eqref{eq:wvoronoi}:
\begin{equation}
  d_w(\mathbf{q}, \mathbf{p}_i)
  = \frac{\lVert \mathbf{q} - \mathbf{p}_i \rVert}{w_i},
  \qquad w_i = \frac{1}{r_i + \varepsilon},
  \label{eq:wvoronoi}
\end{equation}
trong đó $\varepsilon > 0$ tránh chia cho~$0$. Ô Voronoi có trọng số của vật
cản $i$ là tập điểm gần nó nhất theo $d_w$; biên của các ô tạo thành một
mạng đường đi ứng viên có khoảng cách an toàn tương đối đều.

Lớp $\beta$-complexes được dùng để xấp xỉ phần bù rời rạc của các ô, giúp
truy vấn nhanh miền còn trống khi lấy mẫu.

\subsection{Hướng dẫn lấy mẫu và giám sát logic mờ}

RRT* được hướng dẫn bằng cách ưu tiên lấy mẫu trong hành lang an toàn quanh
mạng Voronoi. Bộ giám sát logic mờ nhận ba đầu vào là khoảng cách tới vật
cản gần nhất, tốc độ tương đối và độ cong quỹ đạo, rồi xuất một mức can
thiệp (giữ nguyên, giảm tốc, lập lại kế hoạch). Chi tiết hàm thuộc và luật
mờ được trình bày trong bảng~\ref{tab:rules}.
